\documentclass{article}
\linespread{1.3}
\usepackage[margin=50pt]{geometry}
\usepackage{amsmath, amsthm, amssymb, amsthm, tikz, fancyhdr, float, graphicx, spverbatim, systeme}
\pagestyle{fancy}
\renewcommand{\headrulewidth}{0pt}
\newcommand{\changefont}{\fontsize{15}{15}\selectfont}

\fancypagestyle{firstpageheader}{
  \fancyhead[R]{
    \changefont
    \parbox[t]{4cm}{ % Adjust width as needed
      Michael Huang\\
      EN.555.644.81\\
      Homework N
    }
  }
}

\begin{document}

\thispagestyle{firstpageheader}
{\Large 

\section*{1}
A stock price is currently \$40. It is known that at the end of one month it will be either \$42 or \$38. The risk-free interest rate is 8\% per annum with continuous compounding. What is the value of a one-month European call option with a strike price of \$39?


}
{\Large 

\section*{2}
A stock price is currently \$100. Over each of the next two six-month periods it is expected to go up by 10\% or down by 10\%. The risk-free interest rate is 8% per annum with continuous
compounding. What is the value of a one-year European call option with a strike price of \$100?


}
{\Large 

\section*{3}
For the situation considered in Problem 2, what is the value of a one-year European put option with a strike price of \$100? Verify that the European call and European put prices satisfy put-call parity.


}
{\Large 

\section*{4}
A stock price is currently \$40. It is known that at the end of three months it will be either \$45 or \$35. The risk-free rate of interest with quarterly compounding is 8\% per annum. Calculate the value of a three-month European put option on the stock with an exercise price of \$40. Verify that no-arbitrage arguments and risk-neutral valuation arguments give the same answers.


}
{\Large 

\section*{5}

Consider a European call option on a non-dividend-paying stock where the stock price is \$40, the strike price is \$40, the risk-free rate is 4\% per annum, the volatility is 30\% per annum, and the time to maturity is six months.

\subsection*{(a)} 
Calculate \(u, d\), and \(p\) for a two step tree


\subsection*{(b)}
Value the option using a two step tree.


\subsection*{(c)}
Verify that DerivaGem gives the same answer.


\subsection*{(d)}
Use DerivaGem to value the option with 5, 50, 100, and 500 time steps. 


}

\end{document}